\documentclass[UTF8,a4paper,zihao=-4]{ctexart}
\usepackage[margin=2.2cm]{geometry}
\usepackage{array}
\usepackage{booktabs}
\usepackage{enumitem}
\usepackage{xcolor}
\usepackage{hyperref}
\usepackage{fancyhdr}
\usepackage{lastpage}

\hypersetup{
  colorlinks=true,
  linkcolor=black,
  urlcolor=blue
}

\pagestyle{fancy}
\setlength{\headheight}{15pt}
\fancyhf{}
\lhead{ConsistentFace 论文修改清单}
\rhead{commit cb47cb7}
\cfoot{\thepage/\pageref{LastPage}}

\setlist[itemize]{leftmargin=1.8em,itemsep=0.25em,topsep=0.25em}
\setlist[enumerate]{leftmargin=2.2em,itemsep=0.35em,topsep=0.35em}
\renewcommand{\arraystretch}{1.25}

\title{\textbf{ConsistentFace 论文修改清单}}
\author{基于审稿意见 PDF 与本轮修订生成}
\date{2026 年 7 月 7 日}

\begin{document}
\maketitle

\section*{一、修订范围}

\begin{itemize}
  \item 对照版本：\texttt{3f60883} 到 \texttt{cb47cb7}。
  \item 最新提交：\texttt{cb47cb7 Address review comments and move audit notes}。
  \item 修改文件：\texttt{consistentface\_tip\_zh.tex}。
  \item 新增文件：\texttt{experiment\_audit\_notes.md}。
  \item 变更规模：2 个文件，48 行新增，56 行删除。
\end{itemize}

\section*{二、具体修改清单}

\begin{enumerate}
  \item \textbf{标题重写}
  \begin{itemize}
    \item 位置：\texttt{consistentface\_tip\_zh.tex:33}。
    \item 修改前：标题强调“一种用于双先验人脸恢复的冲突感知适配器”。
    \item 修改后：改为“面向身份保持人脸恢复的参考--文本双先验冲突感知适配器”。
    \item 目的：回应意见中“主题应更明确地指向身份保持、参考图像和文本先验”的要求。
  \end{itemize}

  \item \textbf{摘要整体压缩和重写}
  \begin{itemize}
    \item 位置：\texttt{consistentface\_tip\_zh.tex:45--50}。
    \item 具体修改：把摘要改成三段式结构：问题定义、方法核心、验证结论。
    \item 删除内容：删去过长背景铺垫和过细的失败场景展开。
    \item 新增重点：明确参考图像约束身份、文本提示控制属性，并指出二者失配会导致身份漂移和错误属性采纳。
    \item 目的：回应“摘要应更紧凑、更有力”的意见。
  \end{itemize}

  \item \textbf{贡献列表由四点收敛为三点}
  \begin{itemize}
    \item 位置：\texttt{consistentface\_tip\_zh.tex:71--78}。
    \item 具体修改：将贡献整理为“身份保持双先验冲突建模”“多证据与属性槽路由适配器”“冲突监督和评测协议”三点。
    \item 删除内容：删去贡献项中过多动机性解释，避免贡献列表像背景段落。
    \item 目的：回应“贡献应直接，不要在贡献列表中写动机长段”的意见。
  \end{itemize}

  \item \textbf{相关工作强化最接近工作的差异}
  \begin{itemize}
    \item 位置：\texttt{consistentface\_tip\_zh.tex:92}。
    \item 具体修改：集中补充 CoFaDiff、IConFace、MCS、A$^2$BFR、FaceMe 与本文的差异。
    \item 新增观点：现有方法通常把外部先验作为应强化的条件，而本文把外部先验视为需要验证的证据。
    \item 目的：回应“主 baseline/closest work 要清楚引用，并说明差异”的意见。
  \end{itemize}

  \item \textbf{语言先验不确定性补充依据}
  \begin{itemize}
    \item 位置：\texttt{consistentface\_tip\_zh.tex:97}。
    \item 具体修改：说明文本提示可能来自自动描述、过时人工描述或视觉语言模型生成描述，并可能携带错误属性。
    \item 新增引用方向：结合文本引导恢复和指令编辑文献说明提示错误可能造成身份或属性偏移。
    \item 目的：回应“语言先验不确定性需要文献和任务依据”的意见。
  \end{itemize}

  \item \textbf{多模态可靠性与身份适配器段落重写}
  \begin{itemize}
    \item 位置：\texttt{consistentface\_tip\_zh.tex:100--102}。
    \item 具体修改：补充 CLIP、ControlNet、IP-Adapter、PAIR-Diffusion 等作为多模态融合和条件扩散背景。
    \item 明确差异：IP-Adapter、PhotoMaker、InstantID、PuLID 主要回答如何更强注入身份；本文回答冲突时何时信任参考、何时信任文本、何时回退。
    \item 目的：回应“更强先验”和“可信度分配”区别不够清楚的问题。
  \end{itemize}

  \item \textbf{跨评估器身份指标说明补强}
  \begin{itemize}
    \item 位置：\texttt{consistentface\_tip\_zh.tex:102}。
    \item 具体修改：解释参考 token 使用 ArcFace 特征，若只用 ArcFace 评估会造成输入特征与评价指标耦合。
    \item 新增说明：因此同时报告 AdaFace、MagFace、CurricularFace，并结合属性评估器和人工盲测。
    \item 目的：回应“为什么需要 cross-evaluator identity metrics”的意见。
  \end{itemize}

  \item \textbf{问题表述段落补顺}
  \begin{itemize}
    \item 位置：\texttt{consistentface\_tip\_zh.tex:110--112}。
    \item 具体修改：新增三类决策解释：全局参考/文本偏好、外部先验整体置信度、属性槽局部偏好和回退。
    \item 新增变量说明：$\alpha$ 表示全局参考/文本偏好，$c$ 表示外部先验整体置信度，$\alpha_a,c_a$ 表示属性槽局部偏好和局部回退置信度。
    \item 目的：回应“为什么设计三个动态权重应更顺畅说明”的意见。
  \end{itemize}

  \item \textbf{避免含混术语}
  \begin{itemize}
    \item 位置：\texttt{consistentface\_tip\_zh.tex:112}。
    \item 具体修改：删除“坍缩为单一分数”等容易引发歧义的说法，改为分条说明三类挑战。
    \item 目的：回应“坍缩”等词语表达不够准确的问题。
  \end{itemize}

  \item \textbf{方法开头补充模块来源}
  \begin{itemize}
    \item 位置：\texttt{consistentface\_tip\_zh.tex:123}。
    \item 具体修改：说明 ConsistentFace 基于 SDXL/ControlNet 恢复主干，退化图像由 ControlNet 提供结构约束，参考图像沿用身份编码器和图像提示适配器 token 注入思想，文本由 SDXL 文本编码器提供语义条件。
    \item 明确新增贡献：新增部分不是重新设计去噪主干，而是在条件进入扩散去噪前加入可信度分配适配器。
    \item 目的：回应“方法模块来源不够清楚”的意见。
  \end{itemize}

  \item \textbf{删除正文中的审计占位表}
  \begin{itemize}
    \item 位置：原 \texttt{consistentface\_tip\_zh.tex} 中数据审计表，已删除。
    \item 删除内容：移除包含“待清单导出”“待种子日志补齐”等占位语的数据构成表。
    \item 目的：回应“不要在文中写出 TBD/待清单导出”的要求。
  \end{itemize}

  \item \textbf{主结果表去除 \texttt{TBD}}
  \begin{itemize}
    \item 位置：\texttt{consistentface\_tip\_zh.tex:595--615}。
    \item 具体修改：表 II 的数值改为当前已有均值、$\Delta$ 列、原始 $p$ 值和 Holm 校正 $p$ 值，不再写 \texttt{$\backslash$pm TBD}。
    \item 处理原则：没有真实逐种子标准差时不虚构数值；待补统计移至单独审计说明。
    \item 目的：保持论文正文干净，同时避免伪造未归档统计。
  \end{itemize}

  \item \textbf{删除正文中的缺失日志免责声明}
  \begin{itemize}
    \item 位置：\texttt{consistentface\_tip\_zh.tex:628} 附近。
    \item 具体修改：删除“当前归档缺少逐种子日志，标准差必须补齐”等正文叙述。
    \item 目的：避免正文出现内部审计语气，把投稿前核对事项单独管理。
  \end{itemize}

  \item \textbf{小幅润色结果分析标题句}
  \begin{itemize}
    \item 位置：\texttt{consistentface\_tip\_zh.tex:709}。
    \item 修改前：“这里的关键问题是……”。
    \item 修改后：“本节关注的关键问题是……”。
    \item 目的：去除口语化和批注式表达，使论文语气更正式。
  \end{itemize}

  \item \textbf{新增独立审计说明文件}
  \begin{itemize}
    \item 位置：\texttt{experiment\_audit\_notes.md:1--17}。
    \item 具体内容：单独列出投稿前需要补齐的逐种子统计、评测套件规模、匹配-$s_{RT}$ 反事实套件清单、自然语言/自由文本切片、人工评审材料。
    \item 目的：满足“不要在文中写出 TBD/待清单导出，可以单独生成一个文件阐述”的要求。
  \end{itemize}
\end{enumerate}

\section*{三、正文已清除的占位和内部批注}

\begin{itemize}
  \item 正文 \texttt{consistentface\_tip\_zh.tex} 中未检出：\texttt{TBD}。
  \item 正文 \texttt{consistentface\_tip\_zh.tex} 中未检出：“待清单导出”。
  \item 正文 \texttt{consistentface\_tip\_zh.tex} 中未检出：“待种子日志”“逐种子日志”“当前归档”“审计状态”“标准差需”等内部审计表述。
  \item 审计性内容只保留在 \texttt{experiment\_audit\_notes.md} 中，不进入论文正文。
\end{itemize}

\section*{四、编译和版式检查}

\begin{itemize}
  \item 论文已使用 \texttt{latexmk -pdfxe consistentface\_tip\_zh.tex} 编译成功。
  \item 未发现 LaTeX error、未解析引用、citation warning 或 overfull。
  \item 仅剩字体替换 warning，不影响正文内容完整性。
  \item 最新论文 PDF 已输出到：\texttt{output/pdf/consistentface\_tip\_zh.pdf}。
  \item 已渲染缩略图检查首页、表格页和第 14 页结论/参考文献衔接，未见明显遮挡或大块空白。
\end{itemize}

\section*{五、仍需后续补齐但未写入正文的事项}

\begin{itemize}
  \item 若准备正式投稿，仍需用固定评测清单导出每个随机种子的逐项指标和标准差。
  \item 若补齐后的标准差或配对检验改变显著性结论，应同步更新正文表格、结果分析和结论措辞。
  \item 数据构成表中的身份数、图像数、参考数量分布、冲突类别样本数等应来自真实冻结清单，不应人工补写或虚构。
\end{itemize}

\end{document}
